\section{Predicted Versus Realized Local Readout-Side Changes}
\label{app:causal-validation}

This appendix documents the intervention-validation protocol behind
\Cref{tab:causal-validation-summary}.

\paragraph{Claim under test.}
For a selected contrast direction \(q\) and hidden state \(h\), SRP
asserts that feature \(i\) contributes
\(c_i=\beta_i(q)\,p_i(h)\) to the realized score. If this attribution is
correct, removing the \(d_i\)-component of the hidden state should
change the exact score by \(-c_i\). We therefore ablate
\(h' = h - (h^\top d_i)\,d_i\) for unit-norm decoder direction \(d_i\)
and measure the realized change \(s_\alpha(h')-s_\alpha(h)\) with the
\emph{dense} LM head, so the prediction uses the factorization while the
measurement does not. The claim is falsifiable in three independent
ways: sparse-code misattribution (a contribution assigned to the wrong
feature), decoder non-orthogonality (overlapping features absorbing the
ablated mass), and the row residual (score mass outside the fitted
basis) each break the predicted--realized agreement. Random unit
directions supply the control.

\paragraph{Protocol.}
Per model: approximately 260 contrasts from the main bank; for each, the
top-10 features by absolute contribution plus 10 random-direction
controls, giving 1{,}210--2{,}050 prediction--realization pairs per
model. We report \(r^2\) and the regression slope over pairs passing the
main reading rule (\S\ref{sec:method-reading-local-accounts}), with
bootstrap 95\% CIs; ungated results differ by less than 0.02 in
\(r^2\) throughout. A harness self-test on a residual-free synthetic
factorization recovers \(r^2=1.000\) (slope 1.04) with zero-scoring
controls, so the pipeline itself does not manufacture
agreement.

\begin{table}[!tbp]
\caption{\textbf{Predicted versus realized local readout-side changes.}
For each contrast, the predicted per-feature contribution
\(c_i=\beta_i(q)\,p_i(h)\) is compared with the realized margin change
when \(h\) is ablated along the unit decoder direction \(d_i\)
(readout-side; no re-forward). Per model, \(\sim\)260 bank contrasts
with top-10 features plus 10 random-direction controls each; \(r^2\)
and slope over pairs passing the reading rule, bootstrap 95\% CIs;
Random: control \(r^2\).}
\label{tab:causal-validation-summary}
\centering
\scriptsize
\setlength{\tabcolsep}{3.2pt}
\renewcommand{\arraystretch}{1.06}
\begin{tabular}{@{}lcccc@{}}
\toprule
Model & \(r^2\) [95\% CI] & Slope & Random & \(n\) \\
\midrule
Q-0.8B & 0.886 [.869, .903] & 1.661 & 0.010 & 1610 \\
Q-2B   & 0.834 [.802, .862] & 1.338 & 0.007 & 1820 \\
Q-9B   & 0.829 [.813, .845] & 1.686 & 0.058 & 1860 \\
Min-8B & 0.926 [.913, .937] & 0.957 & 0.009 & 2050 \\
R1-Q-7B & 0.918 [.899, .933] & 1.082 & 0.000 & 1500 \\
R1-L-8B & 0.934 [.920, .947] & 1.126 & 0.034 & 1210 \\
\bottomrule
\end{tabular}
\end{table}

\paragraph{Results and reading.}
All six gated models give \(r^2=0.83\)--\(0.93\) against random-control
\(r^2\le 0.06\) (\Cref{tab:causal-validation-summary}). Slopes are
close to 1 on the 7--8B readouts (0.96--1.13) and rise to 1.34--1.69 on
the Qwen3.5-0.8B/2B/9B readouts: predicted magnitudes there overstate
realized changes while preserving sign and rank. A candidate
explanation is feature interference, with overlapping decoder directions
sharing ablated mass, but we have not isolated it and report slopes as
measured. The scope of the validated claim is deliberately local:
ablations act on the readout side of a realized forward pass (linear in
\(h\), no re-forward), so the result licenses ``feature contributions
predict local readout-side changes,'' not circuit-level necessity,
sufficiency, or generation-level effects
\citep{conmy2023acdc,geiger2023causalabstraction}.
